\subsection{KT-TWOFIELD-FAMILY: an unbounded family of common congruences}
\label{sec:kt-twofield-family}

\begin{proposition}[KT-TWOFIELD-FAMILY]
For every integer $t\geq0$, set
\[
            e_t=268+330t,\qquad w_t=2e_t+1=537+660t.
                                                               \tag{TI.1}
\]
Then, for each $D\in\{3,7\}$,
\[
        \frac{B_{w_t,\chi_D}}{w_t}\in661\mathbf Z_{661}.
                                                               \tag{TI.2}
\]
There is a normalized non-CM primitive newform
$f_{D,t}\in S_{w_t}(\Gamma_0(D),\chi_D)$, with a
coefficient prime $\lambda_{D,t}\mid661$, satisfying
\[
 a_q(f_{D,t})\equiv1+\chi_D(q)q^{w_t-1}
          \pmod{\lambda_{D,t}}
                 \qquad(q\text{ prime},\ q\nmid D).
                                                               \tag{TI.3}
\]
Its normalized three-dimensional classical residual system
$1\oplus(r_{f_{D,t}}\otimes\epsilon^{-e_t})|_{G_{E_D}}$
has semisimplification
$1\oplus\omega^{-e_t}\oplus\omega^{e_t}$.
Consequently no Galois element separates that system from
the Eisenstein triple modulo $661$.
Nevertheless, for every $t\geq0$,
\[
                    K_{536+660t}(\mathbf Z)\{661\}=0.
                                                               \tag{TI.5}
\]
\end{proposition}

\begin{proof}
Use Zhi-Hong Sun, \emph{Congruences concerning Bernoulli
numbers and Bernoulli polynomials}, Discrete Applied
Mathematics 105 (2000), 193--223, Lemma 8.1(a),
pp.~219--220, with $p=661$, character $\chi_D$, and
$n=1$. For any positive $b$ not divisible by $660$ it gives
\[
 (1-\chi_D(661)661^{b-1})\frac{B_{b,\chi_D}}b
 \equiv
 (1-\chi_D(661)661^{b+659})
           \frac{B_{b+660,\chi_D}}{b+660}
                                  \pmod{661}.
                                                               \tag{TI.4}
\]
The conductor is coprime to $661$. Every $b=w_t$ is
positive, exceeds one, and is congruent to $537$ modulo
$660$, so the lemma applies. Both Euler factors in
(TI.4) are $661$-adic units congruent to one. The base
case (TI.2) is KT-TWOFIELD-661. Induction using (TI.4)
proves (TI.2) for all $t$. This argument remains valid
when $661\mid w_t$: the cited congruence concerns the
normalized rational number, so no division by a residue
class of $w_t$ is being performed.

Apply Fretwell--Roberts Theorem 3.6 with
$N=D,M=1,k=w_t,\psi=1,\phi=\chi_D$, and the prime
$(661)$. Its sufficient-condition theorem requires
$k>2$, $661\nmid6D$, and the positive valuation of
$L(1-k,\chi_D)=-B_{k,\chi_D}/k$. These conditions hold
for every $t$. In contrast to their necessary-condition
Theorem 3.5, Theorem 3.6 does \emph{not} impose
$661>k+1$. Primitive nebentypus again makes the resulting
eigenform new. This proves (TI.3).

Since $330\equiv0\pmod3$, every $e_t$ is congruent to
one modulo three. KT-GL2-CM therefore excludes CM at
level three. At level seven, a CM newform would have
zero eigenvalue at the inert prime three. Fermat's theorem
and (TI.3) instead give
\[
 a_3(f_{7,t})\equiv1-3^{536+660t}
                  \equiv1-3^{536}\equiv50\pmod{661}.
\]
The CM classification and inert-prime rule used here
are the same as in KT-TWOFIELD-GLOBAL. Thus both forms
are non-CM. Finally the residual semisimplification follows
from the all-good-primes congruence exactly as in
KT-TWOFIELD-GLOBAL, followed by restriction and twisting.

To prove (TI.5), the already certified residue
$B_{392}\equiv53\pmod{661}$ and the independent short
recurrence certificate $B_{62}\equiv511\pmod{661}$
apply K17-REFLECTOR at $e=268$ and $e=598$, respectively.
Both exponents are even and belong to $[2,658]$.
Thus $A_{661}^{[-268]}=A_{661}^{[-598]}=0$.
The exponent $e_t$ modulo $660$ is $268$ for even $t$
and $598$ for odd $t$. The carried comparison (K2),
which holds for every positive $k$, now gives
$K_{2e_t}(\mathbf Z)\{661\}=0$. This uses periodicity
only of the class-group character index and of the
vanishing criterion, not an isomorphism or order formula
for $K$-groups in different degrees.
\end{proof}

The statement is an infinite-weight obstruction at the
fixed prime $661$. Only its initial member $t=0$ lies in
the large-prime range $661>2e_t+2$ used to define
$\mathcal E_D(e)$ in (TF.8). It does not claim infinitely
many distinct common primes in that range, nor does it
contradict the proposed reflected exclusion (TF.9).
It proves that discarding finitely many weights, changing
the test Galois element, and deleting CM forms cannot
remove this fixed common classical congruence support.


